% Shared exercise chunk: l2-finite-class-consistency
\begin{exercisechunk}
\item \textbf{Generalization by consistency.}
\label[exercise]{ex:finite-class-consistency}
\exerciseprofile{2}{2}{3}{\exproof\quad\excalculation\quad\exinterpretation}{%
Prove the finite-class consistency bound and distinguish a sufficient counting
bound from the sample size a particular learner actually needs.}
Use the assumptions and notation of \cref{thm:finite-class-consistency}.
\begin{enumerate}[(a)]
\item For a predictor $g\in\cF$, compute the probability that $g$ is
consistent with all $n$ observations in terms of $L(g,f)$. Write the
factorization that uses the sampling assumptions.
\item Prove \cref{eq:finite-class-consistency}. You may use the union bound:
$\P(\bigcup_{j=1}^r E_j)\le\sum_{j=1}^r\P(E_j)$ for events
$E_1,\ldots,E_r$. Your proof should apply to every choice among consistent
members, including randomized choices.
\item Compare the guarantees:
\begin{enumerate}[label=(\roman*),leftmargin=*,beginpenalty=10000]
\item Derive the stated sample-size guarantee.
\item Specialize it to the class
of thresholds and compare it with \cref{prop:threshold}. Does this prove
that the threshold learner needs more observations than the sufficient
sample size in \cref{prop:threshold}?
Support your answer with the two bounds.
\end{enumerate}
\end{enumerate}
\end{exercisechunk}
